\chapter{端侧系统集成与实验验证}
\label{cha:verification}

前两章分别围绕双通道脑电采集硬件平台和 EdgeMIFormer 解码算法展开了独立论述：第 \ref{cha:hardware} 章解决了便携式双通道脑电信号的稳定获取、参数配置与无线回传问题，第 \ref{cha:algorithm} 章则回答了标准公开数据集条件下模型结构设计、低通道性能边界以及边缘部署可行性问题。然而，对于面向真实场景的脑机接口系统而言，仅分别完成硬件实现与离线算法验证仍不足以构成完整闭环，更关键的问题在于二者能否在采集通道数、数据格式、实时链路与任务目标上形成一致的系统协同关系。

基于此，本章将研究重点从“单模块可行”进一步推进到“整链路可用”。具体而言，本章首先说明双通道采集终端、BLE 通信链路、上位机软件与解码算法之间的协同逻辑，给出端侧系统的全链路集成方式；随后结合第 \ref{cha:algorithm} 章已经建立的统一实验协议，总结适用于当前系统的实验数据组织方式与评价指标；在此基础上，进一步给出系统实时性与硬件--算法协同实验的验证口径，为后续补充板端时延、真实双导数据与移动场景测试结果建立统一的章节组织框架。

\section{端侧系统全链路集成}

与传统仅针对离线数据集构建的脑电解码流程不同，本文所研究的系统同时包含近体式双通道采集硬件、无线通信链路、配套上位机和面向边缘部署的解码算法，其目标不是单纯提升某一环节的局部指标，而是在有限导联、有限带宽和有限算力约束下维持系统整体可用性。因此，端侧系统集成的核心不在于简单叠加各模块，而在于明确不同层级的职责边界，使硬件能力、数据形态和算法目标相互匹配。

\subsection{硬件-算法协同逻辑}

本文系统采用“采集终端负责稳定获取与传输，解码模型负责任务判别与部署适配”的分层协同思路。双通道采集硬件以 KS1092 与 nRF52832 为核心，重点解决弱脑电信号的前端调理、双通道连续采样、BLE 回传与参数可配置问题；解码算法则以 EdgeMIFormer 为主体，重点解决时空特征提取、低通道条件下的任务适配以及边缘端的轻量化部署问题。两者在系统中并非各自独立存在，而是通过输入导联数、数据帧结构和目标任务定义实现耦合。

这种协同关系首先体现在采集通道与算法目标之间的一致性上。第 \ref{cha:algorithm} 章已经表明，在统一协议下，当输入导联由全通道逐步压缩至 C3/C4 双导时，四分类运动想象任务的性能会出现明显退化，而将任务目标收束为左右手二分类后，低通道条件下的可用性会显著恢复。由此可见，硬件侧“双通道便携化”的设计选择，并不应被简单理解为对标准脑电帽方案的低成本替代，而应与任务目标收束、模型训练策略调整以及系统交互方式的重新定义共同考虑。换言之，端侧系统的协同优化对象不是孤立的模型精度，而是“可观测信息边界”和“实际交互目标”之间的匹配关系。

这种分层协同逻辑可概括为四个层次，如表 \ref{tab:c5_end2end_layers} 所示。采集层关注脑电信号获取与状态感知，传输层关注连续数据回传与可观测性，上位机层承担波形显示、参数配置与实验记录，而解码层则根据离线验证结果选择适合当前硬件约束的任务设置与模型形态。通过这样的层级划分，硬件和算法不再被视为彼此割裂的两项工作，而是共同服务于低通道边缘脑机接口系统的整体目标。

\begin{table}[htbp]
    \centering
    \caption{端侧脑机接口系统全链路中的层级划分与主要职责}
    \label{tab:c5_end2end_layers}
    \begin{tabular}{p{2.8cm}p{4.0cm}p{6.3cm}}
        \hline
        层级 & 主要载体 & 主要职责 \\
        \hline
        采集层 & 双通道脑电硬件终端 & 完成脑电信号接入、前端参数配置、导联状态检测与连续采样 \\
        传输层 & BLE NUS 数据链路 & 完成脑电二进制帧回传、文本命令交互、序列号跟踪与状态可观测 \\
        上位机层 & Qt 上位机程序 & 完成设备连接、帧解析、实时波形显示、频带能量估计、参数下发与实验记录 \\
        解码层 & EdgeMIFormer 及其对照模型 & 完成离线协议验证、低通道任务适配、复杂度分析与后续边缘部署验证 \\
        \hline
    \end{tabular}
\end{table}

需要强调的是，在当前系统阶段，上位机既承担调试工具的职责，也承担实验闭环中的控制与观测职责；而 EdgeMIFormer 的主要定位仍是“面向边缘部署的解码模型主体”，并不要求其在本阶段直接运行于 nRF52832 侧。这样处理的原因在于，采集端主控更适合执行采样、滤波、缓冲与通信等实时任务，而复杂模型推理更适合在 RK3568 等边缘算力节点或上位机侧完成。由此形成的“轻采集端 + 可扩展推理端”结构，更符合当前双通道便携系统的实现条件。

\subsection{数据流转与交互流程}

从系统运行流程看，端侧脑机接口闭环由一条主数据流和一条控制反馈流共同构成。主数据流中，电极采集到的模拟脑电信号首先进入 KS1092 前端完成模拟调理与增益控制，随后由 nRF52832 的片上 ADC 进行双通道采样，经固件侧缓冲区管理与数字滤波后，按照表 \ref{tab:c3_ble_frame} 所示格式打包为脑电二进制帧，并通过 BLE NUS 服务发送给上位机。上位机完成帧解析后，对双通道波形、频带能量与设备状态进行可视化展示，必要时可进一步将数据保存为离线分析文件，用于后续算法验证。

与之对应，控制反馈流由上位机向设备侧发送参数配置与状态查询命令构成。其内容主要包括增益调整、滤波模式切换、缓存状态查询、导联脱落状态回读以及前端复位等操作。由于脑电连续数据流与文本控制流共享同一 BLE 链路，系统通过魔数帧头与文本回执机制完成解析层区分，从而在保持协议轻量化的同时兼顾连续传输和实验控制。结合第 \ref{cha:hardware} 章中的界面设计可知，当前系统已经具备“采集--传输--显示--回控”的完整闭环能力，这也是后续开展低通道真实数据实验的前提。

对于本论文的整体实验组织而言，这一闭环同时支撑两类数据流转模式。第一类是以 BCI Competition IV 2a 为代表的标准公开数据集模式，其作用在于保证算法比较的可重复性与可比性；第二类是以双通道硬件实采数据为代表的系统联调模式，其作用在于验证低通道硬件平台的真实可用性。前者更多服务于算法主体的统一协议评估，后者则更多服务于硬件--算法协同与系统级验证。将这两类模式同时纳入同一章节框架，可以避免把“公开数据集上的算法结论”和“自研硬件上的系统表现”混为一谈。

\section{解码算法性能实验}

在系统集成背景下，算法性能实验不再仅仅回答“某个模型在公开数据集上的最高准确率是多少”，而是需要进一步说明三个问题：其一，离线协议下的比较是否足以支撑模型选择；其二，在低通道硬件约束下，哪些任务设置仍然具有系统可用性；其三，模型复杂度与部署友好度是否与端侧系统的实时性目标相一致。因此，本节在保留第 \ref{cha:algorithm} 章统一协议的基础上，对实验数据组织方式、评价指标与关键结果的系统意义进行归纳。

\subsection{实验数据与评价指标}

当前系统的算法验证主要由两类实验数据构成。一类是标准公开脑电数据集，即第 \ref{cha:algorithm} 章已经采用的 BCI Competition IV 2a 数据集。该数据集具有明确的导联配置、被试划分和评价口径，适合作为不同模型、不同导联数量与不同任务定义之间的统一比较基准。另一类是由第 \ref{cha:hardware} 章所述双通道硬件系统采集得到的真实脑电数据，其主要作用在于支持后续系统级验证，如双导信号可观测性分析、真实链路下的任务可行性验证以及移动场景稳定性测试。

由于两类数据所承担的角色不同，本章对评价指标也做了层次化组织。对于公开数据集上的离线解码实验，重点指标包括分类准确率、Cohen's Kappa 系数、Avg Best Acc 与 Avg Aver Acc，用于评价模型在统一协议下的判别能力和稳定性；对于面向系统集成的部署验证，则进一步关注模型参数规模、量化后模型大小、单次推理延迟以及与链路相关的端到端时延等指标，用于评价模型是否能够在端侧系统中形成可落地的推理方案。表 \ref{tab:c5_metrics} 总结了本章采用的主要评价指标及其对应意义。

\begin{table}[htbp]
    \centering
    \caption{第 5 章实验验证采用的主要评价指标}
    \label{tab:c5_metrics}
    \begin{tabular}{p{3.2cm}p{3.0cm}p{6.7cm}}
        \hline
        指标类别 & 具体指标 & 主要意义 \\
        \hline
        离线解码性能 & Accuracy, Kappa & 衡量不同模型与不同导联配置下的分类判别能力 \\
        训练过程统计 & Avg Best Acc, Avg Aver Acc & 衡量多被试或多轮训练结果的代表性水平与整体稳定性 \\
        模型复杂度 & 参数量、模型大小、MACs/FLOPs & 衡量模型对边缘侧存储与算力资源的占用情况 \\
        推理性能 & 单次推理延迟、量化后加速比 & 衡量模型在边缘推理节点上的实时部署潜力 \\
        系统级指标 & 端到端时延、丢包率、缓冲区溢出次数 & 衡量采集、传输、显示与推理链路的整体闭环表现 \\
        \hline
    \end{tabular}
\end{table}

基于上述分工，本章对于公开数据集实验沿用第 \ref{cha:algorithm} 章表 \ref{tab:chap4_protocol} 所示统一协议，而对于双通道硬件实采数据，则主要将其用于验证采集链路、任务定义与系统稳定性之间的关系，而不在现阶段直接将其与标准公开数据集结果做严格同口径混合比较。这样的安排能够更清楚地区分“算法可比性”与“系统真实性”两种验证目标。

\subsection{与传统算法对比实验}

在统一协议下，第 \ref{cha:algorithm} 章已经对本文模型与低通道条件下的基线模型进行了较为系统的比较。从结果上看，通道退化并不是某一特定模型独有的问题，而是低通道脑电解码任务普遍存在的结构性难点。以统一协议下的 Conformer 型基线为例，当导联由 22 导逐步压缩至 8 导、3 导和 2 导时，四分类准确率呈持续下降趋势，说明空间信息压缩会直接削弱运动想象任务所依赖的判别线索。相应结果与趋势已在图 \ref{fig:degradation_curve_4class} 和表 \ref{tab:channel_degradation} 中给出。

对于端侧系统而言，这组对比结果的意义在于，它并不支持“仅靠模型更换即可完全弥补低通道硬件信息损失”这一假设；更合理的结论是，系统设计必须同时考虑导联配置、任务定义和模型结构三者之间的匹配关系。也正因此，当双导输入保持不变而任务由四分类收束为左右手二分类后，系统可用性能够显著恢复，如图 \ref{fig:c3c4_4class_vs_2class} 和表 \ref{tab:c3c4_task_redefine} 所示。这一结果为第 \ref{cha:hardware} 章所构建的双通道硬件平台指明了更具可行性的算法落点，即优先围绕中央区左右手意图判别构建低通道解码闭环，而不是在硬件约束不变的条件下继续追求高风险的双导四分类任务。

从论文整体结构看，传统算法对比实验的价值不仅在于证明 EdgeMIFormer 是否优于单一基线，更在于说明低通道边缘系统中“什么样的问题值得解”。换言之，实验比较的目标不应仅仅是争取更高的离线分数，而应当服务于系统级任务重定义。这种从系统约束反推算法任务边界的思路，也是本文与单纯算法论文的重要区别。

\subsection{模型轻量化性能分析}

除分类性能外，模型轻量化能力也是端侧系统验证中不可忽视的一环。第 \ref{cha:algorithm} 章已经给出 EdgeMIFormer 在剪枝与量化后的部署导向结果：模型大小由原始尺寸压缩至约 47.9\%，并在 RK3568 NPU 环境下实现了约 31 ms 的单次推理时间。这组结果说明，EdgeMIFormer 在保持较高离线性能的同时，具备进一步进入边缘推理节点的现实可能性。

对于本章而言，这些轻量化结果的意义在于为系统全链路时延分析提供了算法侧边界。也就是说，端侧脑机接口系统的实时性不应仅由 BLE 传输或上位机显示速度决定，解码模型本身的推理耗时同样构成闭环反馈中的关键组成部分。当模型经过量化和剪枝后，其推理耗时已经进入可接受量级，后续系统优化的重点便可自然转移到采样、缓冲、传输与显示等环节之间的协同调度上。

需要指出的是，量化与轻量化结果更多反映的是“模型作为边缘推理模块的部署潜力”，而非整个系统的最终端到端性能。要进一步回答完整闭环是否满足运动想象交互所要求的实时性阈值，还需要将模型推理时间与第 \ref{cha:hardware} 章中的采样周期、BLE 分帧策略、上位机刷新策略结合起来进行联合分析，这也是本章后续系统实时性实验所要关注的重点。

\section{系统实时性实验}

系统实时性实验的目标，是将算法侧已经得到的部署导向结论，与硬件采集链路和上位机显示链路结合起来，形成统一的端到端时延评价框架。与仅统计模型前向推理时间不同，脑机接口系统中的实时性应理解为“从脑电信号产生并被设备采样开始，到结果被显示或可被交互逻辑使用为止”的全流程时延。因此，本节的验证重点不在于单一模块的峰值性能，而在于采样、缓冲、传输、解析和推理之间是否存在明显的瓶颈环节。

\subsection{RK3568 推理延迟测试}

针对 EdgeMIFormer 的端侧部署验证，推理延迟测试应以批大小为 1 的在线推理场景为基本假设，在固定输入长度与固定导联配置条件下，分别统计浮点模型和量化模型在目标推理节点上的单次前向耗时。测试结果建议至少报告平均值、最大值或典型值，以避免仅使用单次最优结果夸大系统实时性。若后续继续采用 RK3568 作为边缘推理节点，则该实验可直接与第 \ref{cha:algorithm} 章的量化结果形成呼应，用于确认模型在实际工具链和目标环境中的稳定表现。

除平均时延外，推理延迟测试还应关注模型输入形态变化对时延的影响。例如，在全通道输入与低通道输入之间，模型复杂度与数据搬运开销可能同时发生变化；而在二分类与四分类设置之间，分类头结构和后处理逻辑也可能影响最终推理耗时。因此，后续板端推理测试不宜仅给出单一配置结果，而应至少选取一个具有代表性的低通道设置，使其与当前双通道硬件系统的目标任务保持一致。

\subsection{全链路延迟优化验证}

相比单纯的模型推理测试，全链路延迟优化验证更能反映端侧脑机接口系统的真实交互能力。该实验应从设备侧采样时间戳出发，依次记录数据写入缓冲区、BLE 分帧发送、上位机接收解析、界面刷新以及解码结果输出等关键节点的时间信息，并据此估计整条链路的累积时延。若条件允许，还可通过示波器、逻辑分析仪或统一软件时间戳机制，对链路中的关键节点进行交叉校验，从而减少单一软件计时方式带来的偏差。

结合第 \ref{cha:hardware} 章已有实现可知，当前系统在固件侧已经具备环形缓冲、样本序列号、时间戳与状态查询机制，在上位机侧也已经具备连续波形显示与状态回读能力，因此后续开展全链路延迟实验具有较好的工程基础。该实验的关键并不是单纯追求某一个局部环节的最小值，而是识别当前闭环中最主要的瓶颈究竟来自采样调度、蓝牙通知、界面刷新还是解码模块本身，并在此基础上决定下一步的系统优化重点。

\section{硬件--算法协同性能实验}

硬件--算法协同性能实验关注的是：在离线公开数据集验证之外，自研双通道硬件系统能否为后续真实意图解码提供足够稳定的输入，以及算法任务定义是否能够与硬件的可观测信息边界形成一致。相比前述离线对比实验，这部分验证更强调“系统能否在真实双导场景中工作”，而非“模型在理想数据集上是否达到最高分数”。

\subsection{两通道真实数据解码验证}

两通道真实数据解码验证宜采用由易到难的组织方式。首先，应利用闭眼/睁眼 Alpha 节律、眨眼伪迹或其他基础生理特征验证采集链路确实能够稳定记录具有神经生理意义的波形变化；在此基础上，再逐步引入与 C3/C4 导联空间位置更匹配的左右手意图判别任务。这样安排的原因在于，真实双导系统的核心风险并不只来自模型结构本身，还来自电极接触质量、运动伪迹、被试状态波动以及无线链路稳定性等多种工程因素。只有先确认输入信号本身具有可解释性，后续的意图解码结果才具有分析意义。

结合第 \ref{cha:algorithm} 章已有结论，真实双导实验的首选任务应优先聚焦左右手二分类，而不宜直接将双导四分类作为主验证目标。这并不是简单的“降低实验难度”，而是让系统实验中的任务设置与采集位置的生理可观测性保持一致。对于当前双通道硬件平台而言，这样的验证路径更有希望得到稳定、可重复且对整篇论文更有支撑力的结果。

\subsection{移动场景稳定性测试}

移动场景稳定性测试的目标，是评估便携式双通道脑电系统在轻度姿态变化、短时移动或连续佩戴条件下的链路连续性和状态稳定性。相较于桌面静态条件下的联调，这类测试更能体现第 \ref{cha:hardware} 章所强调的便携化设计是否真正转化为可用的系统能力。其评价重点包括导联状态是否频繁波动、BLE 连接是否稳定、缓冲区是否出现异常积压，以及上位机界面在较长时段内是否能够保持连续刷新。

从系统验证的角度看，移动场景稳定性并不只是“附加测试项”，而是低通道近体式脑机接口系统区别于传统实验室脑电帽方案的重要证据。若后续实验能够表明系统在轻度移动场景下仍保持稳定采样与连续显示，则不仅能够支撑便携化硬件设计的合理性，也将为进一步开展实时意图交互实验提供必要前提。

\section{本章小结}

本章从系统层面讨论了双通道脑电采集终端、无线通信链路、上位机软件与 EdgeMIFormer 解码算法之间的集成关系，明确了当前端侧脑机接口平台的分层协同逻辑与数据流转方式。在此基础上，本章进一步总结了适用于当前系统的实验数据组织方式与评价指标，并从算法性能、模型轻量化、系统实时性以及硬件--算法协同性能四个方面建立了统一的实验验证框架。

与前几章分别聚焦单一模块不同，本章的核心价值在于将“标准数据集上的离线结论”和“真实双通道硬件上的系统目标”纳入同一分析路径。已有结果表明，低通道场景下的系统设计应优先强调任务目标与采集导联位置的匹配关系，并结合模型轻量化与链路优化共同推进端侧闭环实现。后续随着板端时延、真实双导数据和移动场景测试结果的进一步补充，本章所建立的验证框架可继续扩展为更完整的系统实验章节。
